\documentclass[12pt, a4paper]{article}

% --- Packages ---
\usepackage[margin=2.5cm]{geometry}
\usepackage{booktabs}
\usepackage{array}
\usepackage{tabularx}
\usepackage[table]{xcolor}
\usepackage{amsmath}
\usepackage{graphicx}
\usepackage{hyperref}
\usepackage{fancyhdr}
\usepackage{titlesec}
\usepackage{parskip}
\usepackage{enumitem}
\usepackage{fontenc}
\usepackage{inputenc}
\usepackage[letterspace=150]{microtype}
\usepackage{mdframed}
\usepackage{lmodern}

% --- Colours ---
\definecolor{ehaBlue}{RGB}{0, 84, 142}
\definecolor{ehaLightBlue}{RGB}{224, 237, 248}
\definecolor{ehaGrey}{RGB}{245, 245, 245}
\definecolor{ehaAccent}{RGB}{0, 153, 118}
\definecolor{tableHead}{RGB}{0, 84, 142}

% --- Hyperlink setup ---
\hypersetup{
    colorlinks=true,
    linkcolor=ehaBlue,
    urlcolor=ehaBlue,
    pdftitle={Concept Note: Forecasting Scope and Modelling Approach},
    pdfauthor={Supply Chain Analytics},
}

% --- Header / Footer ---
\pagestyle{fancy}
\fancyhf{}
\fancyhead[L]{\small\textcolor{ehaBlue}{\textbf{EHA Clinics} $\mid$ Procurement Predictive Modelling}}
\fancyhead[R]{\small\textcolor{ehaGrey!60!black}{April 2026}}
\fancyfoot[C]{\small\textcolor{gray}{\thepage}}
\renewcommand{\headrulewidth}{0.4pt}
\renewcommand{\headrule}{\hbox to\headwidth{\color{ehaBlue}\leaders\hrule height \headrulewidth\hfill}}

% --- Section styling ---
\titleformat{\section}
  {\large\bfseries\color{ehaBlue}}
  {\thesection.}{0.6em}{}[\titlerule]

\titleformat{\subsection}
  {\normalsize\bfseries\color{ehaBlue!80!black}}
  {\thesubsection}{0.6em}{}

\titlespacing*{\section}{0pt}{1.4em}{0.6em}
\titlespacing*{\subsection}{0pt}{1em}{0.4em}

% --- Coloured box for key definitions ---
\mdfdefinestyle{keybox}{
  backgroundcolor=ehaLightBlue,
  linecolor=ehaBlue,
  linewidth=1pt,
  leftmargin=0pt,
  rightmargin=0pt,
  innerleftmargin=12pt,
  innerrightmargin=12pt,
  innertopmargin=10pt,
  innerbottommargin=10pt,
  roundcorner=4pt,
}

\mdfdefinestyle{notebox}{
  backgroundcolor=ehaGrey,
  linecolor=ehaAccent,
  linewidth=1.5pt,
  topline=false,
  rightline=false,
  bottomline=false,
  leftmargin=0pt,
  rightmargin=0pt,
  innerleftmargin=12pt,
  innerrightmargin=12pt,
  innertopmargin=8pt,
  innerbottommargin=8pt,
}

% --- Table helpers ---
\newcolumntype{L}[1]{>{\raggedright\arraybackslash}p{#1}}
\newcolumntype{C}[1]{>{\centering\arraybackslash}p{#1}}

\begin{document}

% ============================================================
%  TITLE PAGE
% ============================================================
\begin{titlepage}
  \pagecolor{ehaBlue}
  \color{white}
  \vspace*{2cm}

  {\fontsize{10}{12}\selectfont\MakeUppercase{\textls[100]{Concept Note}}}\\[0.5em]
  {\fontsize{28}{34}\selectfont\bfseries Forecasting Scope\\[0.2em] \& Modelling Approach}\\[1.5em]
  \textcolor{ehaAccent}{\rule{6cm}{2pt}}\\[1.5em]

  {\large Procurement Demand Forecasting Initiative}\\[0.4em]
  {\normalsize EHA Clinics $\cdot$ Supply Chain Analytics}\\[3cm]

  \begin{tabular}{@{}ll}
    \textbf{Document Version} & 1.0 \\[0.4em]
    \textbf{Date}             & April 30, 2026 \\[0.4em]
    \textbf{Phase}            & Q2 Baseline Establishment \\[0.4em]
    \textbf{Status}           & Draft \\
  \end{tabular}

  \vfill
  {\small\textcolor{white!70!ehaBlue}%
    {This document defines the scope boundary for the procurement demand forecasting model,
     including commodity coverage, facility selection, time horizon, forecast granularity,
     and the modelling approach to be applied.}}
\end{titlepage}
\pagecolor{white}\color{black}

% ============================================================
%  CONTENTS
% ============================================================
\newpage
\tableofcontents
\newpage

% ============================================================
%  1. PURPOSE
% ============================================================
\section{Purpose}

This concept note defines the boundaries within which the procurement demand
forecasting model will operate. It answers four foundational questions:
\textit{what} will be forecast, \textit{where}, \textit{over what time period},
and \textit{at what level of detail}.

Establishing a clear scope before modelling begins prevents scope creep,
focuses analytical resources on the data most likely to yield accurate and
actionable forecasts, and provides an agreed reference point for evaluating
whether the KPI target --- a Mean Absolute Percentage Error (MAPE) of
$\leq 20\%$ for critical health commodities by Q3 2026 --- has been met.

% ============================================================
%  2. BACKGROUND
% ============================================================
\section{Background and Context}

\begin{mdframed}[style=keybox]
\textbf{KPI target (as stated by management):}\\
\textit{``Establish current forecast accuracy baseline using historical
Odoo/LoMIS data by Q2, then achieve MAPE $\leq$ 20\% for critical health
commodities across pilot facilities by Q3.''}
\end{mdframed}

\vspace{0.8em}

The primary dataset is the EHA Clinics procurement record extracted from the
Odoo ERP system. It spans \textbf{2018 to 2024}, covers \textbf{16 branch
locations} across Lagos, Abuja, and Kano, and contains more than 50 distinct
procurement categories --- from prescription medicines to vehicle repairs and
staff training.

Because procurement transactions encode both clinical consumption (medicines,
consumables) and institutional overhead (office supplies, maintenance), the
raw dataset cannot be modelled as-is. A well-defined scope is necessary to
isolate the signal that corresponds to patient-driven demand.

% ============================================================
%  3. COMMODITY SCOPE
% ============================================================
\section{Commodity Scope --- What to Forecast}

The forecasting model will target \textbf{critical health commodities}: items
directly consumed in patient care whose demand is driven by consultation
volumes and recurring clinical need. The categories below are identified from
the dataset's \texttt{category\_name} field and are confirmed in-scope for
the pilot.

\vspace{0.8em}
\noindent
\begin{tabularx}{\textwidth}{@{} L{4.2cm} X @{}}
  \toprule
  \rowcolor{tableHead}
  \textcolor{white}{\textbf{Category}} &
  \textcolor{white}{\textbf{Rationale}} \\
  \midrule
  Prescription Medications   & Core clinical demand driver; high-volume, recurring \\
  \addlinespace[2pt]
  Over The Counter Drugs     & High throughput; strong seasonal signal \\
  \addlinespace[2pt]
  Vaccines / NPI Vaccines    & Programme-driven with seasonal and calendar patterns \\
  \addlinespace[2pt]
  Injections                 & Procedure-linked consumable; forecastable against
                               patient-visit data \\
  \addlinespace[2pt]
  Medical Consumables        & Direct patient-care materials (cannulas, gloves, etc.) \\
  \addlinespace[2pt]
  Laboratory Consumables     & Tied to diagnostics throughput; recurring \\
  \addlinespace[2pt]
  Dental Consumables         & Specialty but volume-trackable across pilot sites \\
  \bottomrule
\end{tabularx}

\vspace{1em}

\subsection{Explicitly Excluded Categories}

The following categories are \textbf{out of scope} for forecasting. They are
either non-recurring, driven by administrative or capital decisions rather than
patient demand, or produce so few transactions that modelling is not meaningful:

\begin{itemize}[leftmargin=1.5em, itemsep=2pt]
  \item Office Operations, General Assets \& Expenses, Office Expense
        and Supplies
  \item Vehicle, Facility \& Repairs, Furniture and Fittings
  \item Specialist Services, Outsourced Services, Deliveries
  \item Travel, Rent \& Leases, Events, Research
  \item Services -- Asset Purchase, Laboratory Instruments,
        Medical Equipment (capital items)
  \item \texttt{xxx ARCHIVED xxx} (deprecated catalogue entries)
\end{itemize}

\subsection{Order Status Filter}

Only records where \texttt{state = done} are used. These represent confirmed,
delivered purchase orders that reflect true realised demand. Records with
\texttt{state = purchase} (orders confirmed but not yet receipted) are excluded
from training data to avoid inflating demand estimates.

% ============================================================
%  4. FACILITY SCOPE
% ============================================================
\section{Facility Scope --- Where to Forecast}

The dataset contains 16 named branches. For the pilot phase, forecasting will
focus on the six facilities with the longest continuous procurement history and
sufficient monthly transaction volume to support model training.

\vspace{0.8em}
\noindent
\begin{tabularx}{\textwidth}{@{} L{4.5cm} C{2.5cm} X @{}}
  \toprule
  \rowcolor{tableHead}
  \textcolor{white}{\textbf{Facility}} &
  \textcolor{white}{\textbf{City}} &
  \textcolor{white}{\textbf{Notes}} \\
  \midrule
  Abuja -- Asba and Dantata  & Abuja & Largest transaction volume in the dataset \\
  \addlinespace[2pt]
  Abuja -- Lugbe             & Abuja & Consistent ordering history \\
  \addlinespace[2pt]
  Kano -- Lamido Crescent    & Kano  & High-volume; strong seasonal patterns expected \\
  \addlinespace[2pt]
  Kano -- Independence Road  & Kano  & Representative urban clinic \\
  \addlinespace[2pt]
  Lagos -- Sangotedo Ajah    & Lagos & Only Lagos facility with sufficient history \\
  \addlinespace[2pt]
  REACH Abuja Hub 1 -- Kuje  & Abuja & Community-facing; captures different demand profile \\
  \bottomrule
\end{tabularx}

\vspace{1em}

\begin{mdframed}[style=notebox]
\textbf{Note on Kano Distribution Warehouse:} This branch will be treated as a
\textit{supply-side reference}, not a demand point. Its procurement records
reflect centralised bulk stocking decisions and do not represent direct facility
consumption. It may be used later to model replenishment lead times.
\end{mdframed}

% ============================================================
%  5. TIME HORIZON
% ============================================================
\section{Time Horizon}

\subsection{Training Window}

\textbf{January 2021 -- December 2023} (36 months).

While the dataset begins in 2018, records from 2018--2020 span the clinic
expansion period and the COVID-19 disruption. Including these years without
explicit adjustment would embed abnormal demand patterns --- procurement
surges, facility closures, and supply chain disruptions --- into the training
data. The post-2020 period reflects more stable, representative clinical
operations and is therefore the primary training window.

\subsection{Validation Window}

\textbf{January 2024 -- June 2024} (6 months).

This held-out period is used to compute the Q2 MAPE baseline. Comparing model
predictions against actuals in this window gives an objective, pre-modelling
measure of current forecast accuracy.

\subsection{Forecast Horizon}

Monthly, rolling \textbf{3-month forward} window, extendable to 6 months once
model performance is validated. Monthly granularity aligns with the existing
\texttt{month} column in the dataset and mirrors typical procurement cycle
length at these facilities.

% ============================================================
%  6. FORECAST GRANULARITY
% ============================================================
\section{Forecast Granularity}

The fundamental unit of forecast is:

\begin{mdframed}[style=keybox]
\begin{center}
  \large $\underbrace{\text{Commodity Category}}_{\text{what}}$
  \quad $\times$ \quad
  $\underbrace{\text{Facility}}_{\text{where}}$
  \quad $\times$ \quad
  $\underbrace{\text{Month}}_{\text{when}}$
\end{center}
\end{mdframed}

\vspace{0.8em}

Individual SKU-level forecasting (e.g., \textit{Amoxicillin 500mg $\times$ 100
at Asba -- monthly}) is a natural next step but will not be attempted in the
pilot phase. SKU-level records are currently too sparse at several facilities
to produce statistically stable models, and the measurement overhead would
complicate the Q2 baseline assessment. SKU-level modelling will be scoped
separately once category-level MAPE $\leq 20\%$ is achieved.

% ============================================================
%  7. PERFORMANCE METRIC
% ============================================================
\section{Primary Performance Metric}

Forecast accuracy will be measured using the \textbf{Mean Absolute Percentage
Error (MAPE)}:

\[
  \text{MAPE} = \frac{1}{n} \sum_{t=1}^{n}
  \left| \frac{A_t - F_t}{A_t} \right| \times 100
\]

where $A_t$ is actual procurement quantity in period $t$ and $F_t$ is the
forecast quantity. MAPE is the standard indicator for supply chain forecast
accuracy and is directly aligned with the stated KPI.

MAPE will be computed \textbf{per category-facility pair} to identify which
commodities and locations are under- or over-performing, enabling targeted
model improvements ahead of the Q3 deadline.

% ============================================================
%  8. MODELLING APPROACH
% ============================================================
\section{Modelling Approach}

\subsection{Guiding Principle}

The modelling strategy follows a \textbf{tiered complexity} approach: begin
with the simplest model that is capable of meeting the MAPE $\leq 20\%$
target, and introduce additional complexity only when simpler models
demonstrably fall short. This principle keeps the pipeline interpretable for
supply chain staff, reduces the risk of overfitting on a 36-month training
window, and ensures that model outputs can be explained to --- and trusted by
--- facility managers.

\begin{mdframed}[style=notebox]
\textbf{Core rule:} A more complex model is justified only if it reduces
validation MAPE meaningfully (by at least 3--5 percentage points) compared to
the tier below it. Complexity for its own sake will not be pursued.
\end{mdframed}

\vspace{0.8em}

\subsection{Tier 1 --- Benchmark Models (Baseline)}

Benchmark models are applied first and serve two purposes: they establish the
Q2 MAPE baseline required by the KPI, and they set the performance floor that
all subsequent models must beat to justify their use.

\begin{itemize}[leftmargin=1.5em, itemsep=4pt]
  \item \textbf{Seasonal Na\"{i}ve:} The forecast for month $t$ equals the
        actual value from the same month in the prior year
        ($\hat{F}_t = A_{t-12}$). This is the most widely used operational
        baseline in health supply chain contexts.

  \item \textbf{12-Month Rolling Average:} The forecast equals the mean of
        the preceding 12 months' actuals. Smooths short-term noise without
        requiring any model fitting.
\end{itemize}

If either benchmark already achieves MAPE $\leq 20\%$ for a given
category-facility pair, that pair is considered solved at the baseline stage
and is not escalated to higher tiers.

\subsection{Tier 2 --- Classical Time Series Models}

Classical statistical models are the primary candidates for the pilot phase.
They are well-suited to monthly health commodity data, require relatively
little data volume, and produce uncertainty intervals that are operationally
meaningful for safety-stock planning.

\vspace{0.8em}
\noindent
\begin{tabularx}{\textwidth}{@{} L{3.6cm} X L{3.2cm} @{}}
  \toprule
  \rowcolor{tableHead}
  \textcolor{white}{\textbf{Method}} &
  \textcolor{white}{\textbf{Description and Rationale}} &
  \textcolor{white}{\textbf{Best suited for}} \\
  \midrule

  \textbf{ETS} \newline
  (Error-Trend-Seasonality, \newline Holt-Winters)
  &
  Decomposes each series into error, trend, and seasonal components.
  Parameters are fitted by maximum likelihood; the model automatically
  selects additive or multiplicative seasonality. Well-established for
  monthly pharmaceutical and consumable demand.
  &
  Series with clear trend and/or seasonality; moderate data volume \\

  \addlinespace[4pt]

  \textbf{SARIMA} \newline
  (Seasonal ARIMA)
  &
  Models autocorrelation structures in the series alongside seasonal
  differencing. Captures carry-over demand effects (e.g., a stock-out in
  one month suppressing observed procurement in the next). Order selection
  via AIC/BIC grid search.
  &
  Series with autocorrelation; irregular ordering patterns \\

  \addlinespace[4pt]

  \textbf{Theta Method}
  &
  Decomposes the series into two ``theta lines'' (long-run trend and
  short-run variation) and recombines them. Computationally cheap and
  consistently competitive in empirical benchmarks (M3 and M4
  competitions).
  &
  Sparse or noisy series; rapid iteration \\

  \bottomrule
\end{tabularx}

\vspace{0.6em}

All Tier 2 models will be fitted independently per category-facility pair,
reflecting the expectation that demand dynamics differ meaningfully across
locations (e.g., Kano vs.\ Lagos) and commodity types (e.g., vaccines vs.\
OTC drugs).

\subsection{Tier 3 --- Decomposition and Additive Models}

\textbf{Facebook Prophet} (Meta, open-source) is the designated Tier 3
method. It was designed specifically for business and operational time
series with the following characteristics that closely match this dataset:

\begin{itemize}[leftmargin=1.5em, itemsep=4pt]
  \item \textbf{Irregular seasonality:} Prophet decomposes series into
        yearly, weekly, and custom seasonal components, allowing it to
        capture annual procurement cycles (e.g., rainy-season disease
        burden, year-end budget spending) without manual specification.

  \item \textbf{Holiday and event effects:} Nigerian public holidays
        (Eid al-Fitr, Eid al-Adha, Christmas, Independence Day) and
        facility-specific events (donor supply arrivals, audit periods) can
        be encoded as regressors, preventing systematic forecast bias around
        these dates.

  \item \textbf{Missing data and outliers:} Prophet is robust to gaps in
        the time series --- relevant here because some facility-month
        combinations have zero or missing procurement records due to
        stockouts or reporting lapses rather than true zero demand.

  \item \textbf{Uncertainty intervals:} Prophet natively produces
        credible intervals around each forecast, which can be used directly
        to set reorder point buffers.
\end{itemize}

Prophet will be applied to category-facility pairs where Tier 2 models fail
to achieve MAPE $\leq 20\%$, and as a parallel track to validate Tier 2
results on the highest-volume commodity categories.

\subsection{Tier 4 --- Machine Learning Models (Conditional)}

Machine learning (ML) models will be evaluated \textbf{only if} the
three tiers above do not meet the MAPE target across the majority of
category-facility pairs. The primary candidates are:

\begin{itemize}[leftmargin=1.5em, itemsep=4pt]
  \item \textbf{Gradient Boosting (XGBoost / LightGBM):} Uses lag
        features (e.g., $A_{t-1}, A_{t-3}, A_{t-12}$), rolling
        statistics, and categorical covariates (facility, commodity
        category, quarter) as predictors. Can model cross-facility
        patterns jointly, which is an advantage when individual facility
        series are short. Requires careful feature engineering and
        time-respecting cross-validation to avoid data leakage.

  \item \textbf{Random Forest Regression:} Similar feature structure to
        gradient boosting but less prone to overfitting on small datasets.
        Serves as a lower-variance ML alternative.
\end{itemize}

\begin{mdframed}[style=notebox]
\textbf{Why deep learning is deferred:} Recurrent neural networks (LSTM,
Temporal Fusion Transformers) require substantially more data per series
than is currently available (36 months at category level). They also offer
limited interpretability for operational decision-making. These approaches
will be reconsidered at the SKU-level phase when data volume is higher.
\end{mdframed}

\subsection{Model Selection and Validation Strategy}

Model selection will use \textbf{walk-forward (expanding-window)
cross-validation} on the training data (2021--2023), followed by final
evaluation on the held-out validation window (Jan--Jun 2024).

The procedure for each category-facility series is:

\begin{enumerate}[leftmargin=1.5em, itemsep=4pt]
  \item Train on Jan 2021 -- Dec 2022; forecast Jan--Jun 2023; compute MAPE.
  \item Extend training to Jun 2023; forecast Jul--Dec 2023; compute MAPE.
  \item Average the two cross-validation MAPEs to rank candidate models.
  \item Retrain the best-ranked model on the full training window
        (Jan 2021 -- Dec 2023) and evaluate on the validation window
        (Jan--Jun 2024) to report the official Q2 baseline.
\end{enumerate}

The winning model may differ by category-facility pair. A lookup table
mapping each pair to its selected model will be maintained and updated as
part of ongoing monitoring.

\vspace{0.6em}
\noindent
\begin{tabularx}{\textwidth}{@{} L{3.8cm} X @{}}
  \toprule
  \rowcolor{tableHead}
  \textcolor{white}{\textbf{Tier}} &
  \textcolor{white}{\textbf{Models}} \\
  \midrule
  Tier 1 -- Benchmark      & Seasonal Na\"{i}ve, 12-Month Rolling Average \\
  \addlinespace[2pt]
  Tier 2 -- Classical      & ETS (Holt-Winters), SARIMA, Theta \\
  \addlinespace[2pt]
  Tier 3 -- Additive       & Facebook Prophet (with holiday regressors) \\
  \addlinespace[2pt]
  Tier 4 -- ML (if needed) & XGBoost, LightGBM, Random Forest \\
  \addlinespace[2pt]
  Deferred                 & LSTM, Temporal Fusion Transformers \\
  \bottomrule
\end{tabularx}

% ============================================================
%  9. SCOPE SUMMARY
% ============================================================
\section{Scope Summary}

\vspace{0.4em}
\noindent
\begin{tabularx}{\textwidth}{@{} L{3.8cm} X @{}}
  \toprule
  \rowcolor{tableHead}
  \textcolor{white}{\textbf{Dimension}} &
  \textcolor{white}{\textbf{Decision}} \\
  \midrule
  Commodities       & 7 health commodity categories (see Section 3) \\
  \addlinespace[2pt]
  Facilities        & 6 pilot facilities across Abuja, Kano, Lagos \\
  \addlinespace[2pt]
  Training data     & January 2021 -- December 2023 \\
  \addlinespace[2pt]
  Validation data   & January 2024 -- June 2024 \\
  \addlinespace[2pt]
  Forecast horizon  & Monthly; 3-month rolling window \\
  \addlinespace[2pt]
  Granularity       & Category $\times$ Facility $\times$ Month \\
  \addlinespace[2pt]
  Order filter      & \texttt{state = done} only \\
  \addlinespace[2pt]
  Accuracy metric   & MAPE per category-facility pair \\
  \addlinespace[2pt]
  Modelling tiers   & Seasonal Na\"{i}ve $\to$ ETS/SARIMA/Theta $\to$
                      Prophet $\to$ ML (escalate as needed) \\
  \addlinespace[2pt]
  Validation method & Walk-forward cross-validation; final eval on
                      Jan--Jun 2024 hold-out \\
  \bottomrule
\end{tabularx}

% ============================================================
%  9. OPEN QUESTIONS
% ============================================================
\section{Open Questions for Stakeholder Alignment}

The following questions require input from clinical and supply chain
stakeholders before this scope is finalised:

\begin{enumerate}[leftmargin=1.5em, itemsep=6pt]
  \item \textbf{REACH Hub inclusion:} Are the REACH Hub clinics part of the
        formal pilot programme, or do they operate under a separate
        procurement protocol that would require a distinct forecasting track?

  \item \textbf{Kano Distribution Warehouse:} Does the warehouse dataset
        reflect what was \textit{dispatched to facilities} (downstream demand)
        or what was \textit{purchased from suppliers} (upstream procurement)?
        The answer determines whether it can serve as a proxy for facility-level
        consumption.

  \item \textbf{Critical commodity list:} Does the clinical or supply chain
        team maintain a formal list of ``essential'' or ``critical'' health
        commodities (analogous to the WHO Essential Medicines List or a
        national formulary)? If so, this list should anchor the in-scope
        product selection rather than category-level filtering alone.
\end{enumerate}

% ============================================================
%  10. NEXT STEPS
% ============================================================
\section{Next Steps}

Upon sign-off of this scope document, the following activities will commence:

\begin{enumerate}[leftmargin=1.5em, itemsep=6pt]
  \item \textbf{Data extraction and cleaning} --- Filter the procurement
        dataset to the defined scope; handle missing values, zero-quantity
        records, and duplicate entries.
  \item \textbf{Exploratory data analysis} --- Characterise demand
        distributions, identify seasonality patterns, and flag facilities or
        categories with insufficient history.
  \item \textbf{Baseline MAPE computation} --- Apply a na\"{i}ve benchmark
        model (e.g., prior-year same-month average) to the validation window
        and record the Q2 baseline MAPE per category-facility pair.
  \item \textbf{Model training and selection} --- Apply the four-tier
        modelling framework (Section 8) using walk-forward cross-validation;
        record the winning model per category-facility pair and report the
        official Q2 baseline MAPE.
  \item \textbf{Monitoring pipeline} --- Set up a repeatable monthly
        re-evaluation process to track MAPE improvement toward the Q3
        target of $\leq 20\%$.
\end{enumerate}

\vfill

\noindent\textcolor{ehaBlue!40!white}{\rule{\textwidth}{0.4pt}}\\[0.4em]
{\small\textcolor{gray}{EHA Clinics $\cdot$ Procurement Predictive Modelling $\cdot$
Concept Note v1.0 $\cdot$ April 2026}}

\end{document}
